\begin{sgastatement}{Corollary 2.10}
\(\pi\) is ``generically associative'' and \(\sigma\) is a ``generic
inverse'': if \(x,y,z\in\Omega(S)\), and if the expressions below are
defined (as they always are over an open subset of \(\Omega\) containing
the unit section), then
\[
 \pi(x,\pi(y,z))=\pi(\pi(x,y),z),\qquad
 \pi(x,\sigma(x))=e=\pi(\sigma(x),x).
\]
\end{sgastatement}
\SGARefTarget{sga3:xxv:corollary:2:10}

Indeed, the two sides of each formula define morphisms between smooth
\(\mathbb Z\)-schemes of finite presentation, and the morphisms agree on
the generic fibers by \SGARefLink{sga3:xxv:proposition:2-9:item:iii}{2.9(iii)}.\footnote{The French source reads
``by \SGARefLink{sga3:xxv:proposition:2-9:item:iii}{2.10(iii)}''; the intended reference is Proposition 2.9(iii).}

\begin{sgastatement}{\SGAInlineRefTarget{sga3:xxv:corollary:2-11}Corollary 2.11}
Let \(\alpha\in R\), and let \(S\) be a scheme. Suppose that
\(x,y\in\mathbb G_a(S)\) satisfy
\[
 (p_\alpha(x),p_{-\alpha}(y))\in V_1(S),
 \qquad 1+xy\in\mathbb G_m(S).
\]
These conditions define an open subset of \(\mathbb G_{a,S}^2\)
containing the section \((0,0)\). Then
\[
 \pi(p_\alpha(x),p_{-\alpha}(y))
 =
 \left(
 p_{-\alpha}\!\left(\frac{y}{1+xy}\right),
 \alpha^\ast(1+xy),
 p_\alpha\!\left(\frac{x}{1+xy}\right)
 \right).
\]
\end{sgastatement}
\SGARefTarget{sga3:xxv:corollary:2:11}

The proof is the same as above, by \SGARefLink{sga3:xx:2.1}{Exposé XX, 2.1}.

\SGASection{3}{Existence theorem: conclusion of the proof}
\SGARefTarget{sga3:xxv:section:3}

For convenience, make the following definition.

\begin{sgastatement}{\SGAInlineRefTarget{sga3:xxv:definition:3-1}Definition 3.1}
Let \(S\) be a scheme and \(G\) an \(S\)-group scheme. The group \(G\)
is called \emph{admissible} if there is an open immersion of
\(S\)-schemes
\[
 i\colon\Omega_S=\Omega\times S\longrightarrow G
\]
such that:
\begin{enumerate}[label=\textup{(\roman*)},leftmargin=2.2em]
\item the diagram
\[
\begin{tikzcd}[column sep=4.5em,row sep=3.2em]
(V_1)_S
  \arrow[r]
  \arrow[d,"\pi"]
&
\Omega_S\times_S\Omega_S
  \arrow[r,"i\times_S i"]
&
G\times_SG
  \arrow[d,"\pi_G"]
\\
\Omega_S
  \arrow[rr,"i_S"]
&&
G
\end{tikzcd}
\]
is commutative, where \(\pi_G\) denotes the multiplication morphism
of \(G\);
\item there is a finite set of sections \(a_j\in\Omega(S)\) such that
the open subsets \(i(a_j)\,i(\Omega_S)\) cover \(G\).
\end{enumerate}
\end{sgastatement}
\SGARefTarget{sga3:xxv:definition:3:1}

By ``Weil's theorem'' (\SGARefLink{sga3:xviii:corollary:3-13:item:iii}{Exposé XVIII, 3.13(iii)} and (iv)), we have:

\begin{sgastatement}{Lemma 3.2}
If every admissible \(S\)-group is affine whenever \(S\) is étale and
of finite type over \(\mathbb Z\), then there exists an admissible
affine \(\mathbb Z\)-group.
\end{sgastatement}
\SGARefTarget{sga3:xxv:lemma:3:2}

We now have:

\begin{sgastatement}{Lemma 3.3}
Let \(S\) be a scheme and \(G\) an admissible \(S\)-group. Then \(G\)
is smooth and of finite presentation over \(S\), with affine,
connected, semisimple fibers.
\end{sgastatement}
\SGARefTarget{sga3:xxv:lemma:3:3}

Since \(\Omega_S\) is smooth and of finite presentation over \(S\),
with connected fibers, the same is true of \(G\) by condition (ii).
To prove \SGARefLink{sga3:xxv:lemma:3:3}{3.3}, we may therefore suppose that \(S\) is the spectrum of a
field \(K\). Identify \(\Omega_K\) with its image in \(G\). Plainly
\(\Omega_K\) is the product
\[
 \left(\prod_{\alpha\in R^-}U_{\alpha,K}\right)
 T_K
 \left(\prod_{\alpha\in R^+}U_{\alpha,K}\right)
\]
of the subgroups \(T_K\) and \(U_{\alpha,K}\), \(\alpha\in R\), of
\(G\). Thus the Lie algebra of \(G\) is identified with the direct sum
\[
 \operatorname{Lie}(T_K)
 \oplus
 \bigoplus_{\alpha\in R}\operatorname{Lie}(U_{\alpha,K}).
\]
The inner automorphism defined by a section of \(T_K\) acts on
\(U_{\alpha,K}\), and hence on \(\operatorname{Lie}(U_{\alpha,K})\),
through the character
\[
 \alpha\in R\subset M
 \simeq\operatorname{Hom}_{K\text{-gr.}}(T_K,\mathbb G_{m,K}).
\]
The preceding decomposition of \(\operatorname{Lie}(G)\) is therefore
exactly its decomposition under the adjoint action of \(T\). Thus the
roots of \(G_K\) with respect to \(T_K\) are the elements
\(\alpha\in R\).

Apply \SGARefLink{sga3:xix:corollary:1-13}{Exposé XIX, 1.13}. Let \(T_\alpha\) be the maximal torus of
\(\ker(\alpha)\subset T_K\), and put
\[
 Z_\alpha=\operatorname{Centr}_G(T_\alpha).
\]
It is enough to prove that each \(Z_\alpha\) is reductive. But
\[
 Z_\alpha\cap\Omega_K
 =
 \left(
 \prod_{\substack{\beta\in R^-\\\beta|_{T_\alpha}=e}}
 U_{\beta,K}
 \right)
 T_K
 \left(
 \prod_{\substack{\beta\in R^+\\\beta|_{T_\alpha}=e}}
 U_{\beta,K}
 \right).
\]
The roots vanishing on \(T_\alpha\) are the rational multiples of
\(\alpha\), hence precisely \(\alpha\) and \(-\alpha\). Consequently,
\[
 Z_\alpha\cap\Omega_K
 =U_{-\alpha,K}\,T_K\,U_{\alpha,K}.
\]
By \SGARefLink{sga3:xx:3.4}{Exposé XX, 3.4}, to prove \(Z_\alpha\) reductive it is enough to show
that \(U_{\alpha,K}\) and \(U_{-\alpha,K}\) do not commute, and this
follows immediately from \SGARefLink{sga3:xxv:corollary:2-11}{2.11}.

It follows from \SGARefLink{sga3:xxv:lemma:3:2}{3.2} and \SGARefLink{sga3:xxv:lemma:3:3}{3.3} that the proof will be complete once the
following lemma is proved.

\begin{sgastatement}{Lemma 3.4}
Let \(S\) be a locally noetherian scheme of dimension at most \(1\),
and let \(G\) be a smooth \(S\)-group of finite type, with affine,
connected, semisimple fibers. Then \(G\) is affine (and hence
semisimple).
\end{sgastatement}
\SGARefTarget{sga3:xxv:lemma:3:4}

\noindent\textit{Nota.}\quad
In Exposé XVI it was shown that \SGARefLink{sga3:xxv:lemma:3:4}{3.4} holds without any hypothesis on
\(S\), but the proof is comparatively delicate. Since only the special
case \SGARefLink{sga3:xxv:lemma:3:4}{3.4} is needed here, we give a direct proof.

Consider the Lie algebra \(\mathfrak g\) of \(G\), a locally free
\(\mathcal O_S\)-module, and the adjoint representation
\[
 \operatorname{Ad}\colon
 G\longrightarrow\operatorname{GL}_{\mathcal O_S}(\mathfrak g).
\]
To prove that \(G\) is affine over \(S\), it is enough to prove that
\(\operatorname{Ad}\) is affine. Since \(G\) is smooth with connected
fibers, it is separated over \(S\) (Exposé VI\textsubscript{B}, \SGARefLink{cor:sga3-VIB-5.5}{5.5}),
so \(\operatorname{Ad}\) is separated. By the result proved in the
Appendix (see \SGARefLink{sga3:xxv:proposition:4:1}{4.1}), it suffices to show that \(\operatorname{Ad}\) is
quasi-finite. We are thus reduced to the case in which \(S\) is the
spectrum of a field. In that case \(G\) is affine and therefore
semisimple, and the assertion is \SGARefLink{sga3:xxii:prop:5.7.14}{Exposé XXII, 5.7.14}.
